\subsection{Finite fusion with the named Klein-four labels retained}

\begin{result}{A.4.35}{Proved}{An explicit finite refinement recognizes every labelled Klein-four lift at a prescribed lower quotient}
Use the actual groups $G_t$, their named $H=C_2^2$, and
the fixed-pair forgetting maps of A.4.34.
For every specified finite quotient $q:G_t\to Q$ there
is a specified finite refinement $q^\dagger:G_t\to Q^\dagger$
with the following property.  Retain the actual label
permutation map in the refinement.  Every homomorphism
\[
 \iota:H\longrightarrow Q^\dagger
\]
whose elementwise label permutations are the named
permutations of $H$ has an actual $g\in G_t$ such that
\[
 q(g)\,\operatorname{pr}_Q(\iota(h))\,q(g)^{-1}=q(h)
                    \qquad(h\in H).                    \tag{KFU.1}
\]
The input groups, their quotient maps and finite bounds
are constructed below.  At $t=3$, the correcting element
can be chosen as a pure mapping class followed or preceded
by at most the fixed half twist interchanging the two
labels of $C_0$.  This last half twist is killed by every
four-point forgetting map used in (A.T17)--(A.T20).
The assertion concerns the specified finite refinement;
it does not assert conjugacy at the original output
without that retained input.
\end{result}

\begin{proof}
\par\medskip\noindent\textbf{The one-point torsion comparison.}\quad
The same finite construction as A.4.33 has the following
additional property.  Let $G_p$, $F$, and $a$ be as in
(BCF.2)--(BCF.5), and let $u:G_p\to U$ be the refinement
of a given $q_0:G_p\to Q_0$.  Every element $x\in U$
of order dividing two that belongs to $u(Fa)$ has lower
image in $Q_0$ conjugate by an element of $q_0(F)$ to
the image of an actual order-two vertex element of
\[
 E=F\rtimes\langle a\rangle
       \cong C_2*F_s\quad\hbox{or}\quad C_2*C_2*F_s.
                                                               \tag{KFU.2}
\]
The number of vertex types is respectively one or two.
The construction of $u$ is uniform in the named ordinary
half-turn fixing $p$: its characteristic kernels depend
only on $F$ and the input finite quotient.

To prove the assertion, choose an actual representative
$fa$ of $x$.  Since its action on $W=F/J_2$ has square
one, $(fa)^2\in F$ acts trivially by conjugation on $W$.
It therefore lies in $Z(W)$ at this level.  The center
comparison (BCF.7) gives $(fa)^2\in J_1$.
Thus $fa$ has order two in $E/J_1$; its image in $E/J$
is nonidentity, because it maps to the nonidentity of
$E/F=C_2$.  The order-two assertion of A.4.7 now places
that image in the conjugacy class of one of the actual
vertex involutions in $E/J$.  This projects to $Q_0$
because $J\subseteq\ker(q_0|_F)$.
The conjugator can be chosen in $F$: for any vertex
involution $a_i$, the equality $E=F\langle a_i\rangle$
lets one remove the $a_i$ factor from a conjugator
without changing its conjugation of $a_i$.
This proves (KFU.2), including the actual source of
the conjugator and its lower quotient map.

\par\medskip\noindent\textbf{Torsion recognition over an unordered fixed pair.}\quad
Fix one step $\pi_t:G_t\to G_{t-1}$, with new pair
$\{p,q\}$ and involution $a=a_t$.
For every specified output $v:G_t\to V$, we now construct
a finite refinement $w:G_t\to W$ such that every
order-two element $x\in W$ with the following two
properties has lower image conjugate to $v(a)$ by an
actual element of $K_t$:
\[
 \text{$x$ has the named label permutation of $a$, and}
 \quad x\in w(K_ta).                                    \tag{KFU.3}
\]
Retain the label map in $v$ from the start.
Let $G_t^0$, $G_q$, $F_p$, and $F_q$ be the actual groups
and maps in (HCF.7).
Apply the one-point construction (BCF.2)--(BCF.3) to
$v|_{G_t^0}$ and $F_p$, obtaining $u_1:G_t^0\to U_1$.
The quotient $\overline U_1=U_1/u_1(F_p)$ is a specified
finite quotient of $G_q$.
Apply the same construction to $\overline U_1$ and $F_q$,
obtaining $u_2:G_q\to U_2$.
Set
\[
 \overline U_2=U_2/u_2(F_q),\qquad
 T=\operatorname{im}(u_1,u_2\pi_p).
                                                               \tag{KFU.4}
\]
The first group is a specified finite quotient of
$G_{t-1}$.  The second is an actual finite image of
$G_t^0$.  Induce its homomorphism to $G_t$ using the
index-two cosets, and define $W$ to be the actual joint
image of
\[
       v,\qquad \operatorname{Ind}_{G_t^0}^{G_t}(T),
                    \qquad \overline U_2\pi_t.          \tag{KFU.5}
\]
In particular all three maps, including the last lower
map, have been explicitly retained.

An element $x$ as in (KFU.3) has an actual representative
in $G_t^0$, because its label permutation fixes $p,q$.
Its $U_2$ coordinate has order two and belongs to the
coset $u_2(F_q a)$, by its $\overline U_2$ coordinate.
At this one-point stage both fixed points of $a$ were
unmarked on the old sphere.  Thus (KFU.2) gives exactly
two possible ordinary vertex lifts.  Choose once an
actual half twist $P\in K_t$ along an arc joining $p$
and $q$.  The two lifts can be taken as the images of
\[
                         a,\qquad PaP^{-1}             \tag{KFU.6}
\]
under forgetting $p$.  This identification includes the
base paths: the path of $q$ under the half twist joins
one of the two fixed points to the other, and is the
base path to the second cone-point generator in the
orbifold presentation (KFU.2).  The two fixed vertices
therefore give precisely the two elements in (KFU.6),
with no discarded conjugation; any changed path is
an actual $F_q$ conjugation already allowed by (KFU.2).

The resulting $F_q$ conjugator lifts to an actual element
of $K_t^0$, by the surjection $K_t^0\to F_q$ in the
ordered configuration fibration.  Conjugate $x$ by that
lift.  Its image in $\overline U_1$ is now the image of
one of the two actual elements (KFU.6).  Hence its
$U_1$ coordinate lies in the corresponding coset of
$u_1(F_p)$.  It still has order two.
Apply (KFU.2) once more, now to the point $p$.
The other fixed point $q$ is marked and deleted from
the point-pushing surface, so (KFU.2) has exactly one
order-two vertex type.  The lower image in $V$ is
therefore conjugate, by an actual element of $F_p$,
to the selected element in (KFU.6).
In the second case conjugate once more by $P^{-1}$.
All these conjugators are actual elements of $K_t$.
This proves (KFU.3)--(KFU.5) with the fixed-vertex switch
fully retained.

\par\medskip\noindent\textbf{The four-point starting group.}\quad
By (HCF.4), $G_0=F_2\times H$ with the actual deck
factor $H$.  For a finite quotient $q$ put $P=q(F_2)$
and $L=\ker(F_2\to P)$.  The specified finite input
\[
                    (F_2/L'L^2)\times H               \tag{KFU.7}
\]
refines $q$: its map to $Q$ sends $(\bar f,h)$ to
$q(f)q(h)$.  Every order-two element of $F_2/L'L^2$
has trivial image in $P$.  Here is the required graph
torsion proof.  An element whose image generates a
nontrivial $C_2$ in $P$ would lie over the free subgroup
$F_*\leq F_2$ which is the inverse image of that $C_2$.
The map $F_*\to C_2$ lifts to $\mathbb Z/4$, by
choosing lifts of the images of a free basis.  It kills
$L'$ and $L^2$, since $L$ maps into $2\mathbb Z/4$.
It would therefore send the alleged order-two element
to an odd residue of order four, a contradiction.
Thus every labelled homomorphism $H\to$(KFU.7) projects
pointwise to the actual $H$ in $Q$.  No conjugator is
needed in the starting case.  If $|P|=m_0$, the input
order is $4m_0 2^{m_0+1}$, since $L$ has free rank
$m_0+1$.

\par\medskip\noindent\textbf{The complete fusion recursion.}\quad
For the given $q:G_t\to Q$, first retain its label map
and form the unordered-pair fibre detector
$v:G_t\to V$ of (HCF.6), relative to that output.
Thus
\[
 \operatorname{im}\bigl(C_V(v(a_t))\cap v(K_t)
                                      \longrightarrow Q\bigr)=1.
                                                               \tag{KFU.8}
\]
Apply (KFU.4)--(KFU.5) to this $v$ and call the resulting
finite quotient $w:G_t\to W$.  Form the actual lower
quotient $\overline W=W/w(K_t)$ of $G_{t-1}$.
Use the already constructed fusion refinement at level
$t-1$ for that quotient.  The final input is the actual
joint image of $w$ and this refined lower map composed
with $\pi_t$:
\[
 G_t\longrightarrow
 W\times\mathcal F_{t-1}(\overline W).
                                                               \tag{KFU.9}
\]
This defines $\mathcal F_t(q)$ recursively, with (KFU.7)
as its starting value.

Let $\iota$ be a homomorphism in the statement at this
input.  Its lower projection has the named elementwise
permutations on the remaining marks, and remains
faithful because those on $D_0$ already are faithful.
The induction hypothesis gives an actual element of
$G_{t-1}$ which conjugates its $\overline W$ projection
to the named $H$.  Lift that element through the actual
surjection $G_t\to G_{t-1}$ and make this conjugation.
Every permutation in $G_t$ centralizes the named label
action of $H$, so the prescribed elementwise labels
are unchanged.
The $W$ projection of $\iota(a_t)$ now lies in the
coset $w(K_ta_t)$ and has the named labels.  The pair
torsion comparison provides an actual $k\in K_t$
which conjugates its $V$ projection to $v(a_t)$.
Conjugate the full homomorphism by this $k$.
Choose either of the other two generators $b_t\in H$.
Its $V$ projection commutes with $v(a_t)$ and has the
same projection as $v(b_t)$ in $V/v(K_t)$, because
the $K_t$ correction did not change the forgotten
base.  Therefore
\[
           \operatorname{pr}_V(\iota(b_t))v(b_t)^{-1}
                     \in C_V(v(a_t))\cap v(K_t).       \tag{KFU.10}
\]
Formula (KFU.8) kills its image in $Q$.
Both generators are consequently the named generators
in $Q$.  Their product is as well.  This proves (KFU.1)
by induction, with actual ordinary conjugators at
every step.  In particular the second generator was
compared in the retained centralizer fibre (KFU.10),
not merely under an unrestricted Birman projection.

\par\medskip\noindent\textbf{Labels of the final correcting element.}\quad
At $t=3$, the permutation image of $N_3$ has order
sixteen in the order-thirty-two permutation centralizer.
The exact discriminant calculation in A.1 shows that
it consists of the deck translations on $D_0$ together
with the even-parity flips of the three fixed pairs:
loops around the two ordinary parameter punctures
flip $(A_0,C_0)$ and $(B_0,C_0)$.
Consequently any permutation in $G_3$ differs from
the identity or the single $C_0$ flip by the permutation
of an actual element of $N_3$.
Multiply the correcting element on the left by such
an element of $N_3$.  This preserves (KFU.1), because
$N_3$ centralizes the named $H$, and leaves a pure
permutation or the $C_0$ flip.  Removing the chosen
actual half twist $P_C$ gives a pure element, so the
corrector has the asserted form $uP_C^\epsilon$ or
$P_C^\epsilon u$, with $u$ pure and $\epsilon\in\{0,1\}$.
The support disk of $P_C$ contains only the two labels
of $C_0$.  Deleting that pair kills this actual mapping
class, proving the final compatibility assertion.

\par\medskip\noindent\textbf{Finite instructions and bounds.}\quad
All operations in (KFU.4)--(KFU.9) are characteristic
free-group homology quotients, finite induced maps,
actual image joins, or quotients by the image of a
specified normal point-pushing kernel.
For an input of order $m$ at a pair stage, the one-point
bounds are those of (HCF.11).  If $m_1\leq B(m,2t+2)$
and $m_2\leq B(m_1,2t+1)$, (KFU.5) has order at most
\[
                         2m\,m_1^2m_2^3.              \tag{KFU.11}
\]
First substitute the bound (HCF.12) for the fibre input
$m$; then multiply this bound by the recursively bounded
input at $t-1$, as in (KFU.9).
The starting bound is at most $4m2^{m+1}$.
There are three pair stages.  This supplies an explicit
finite bound and an exact source-map construction, rather
than a compactness selection of an unknown level.
It does not claim that the remaining finite pentagon
shadow enumeration has been performed.
\end{proof}
